\cotefile{10}
\scanpage{185}

% slanted
\margin{$k$ \unsure{infini} ???}

Soit $k$ un corps. $k$ est un groupe algébrique sur $k$, dont nous cherchons les représentations linéaires, i.e les applications \uline{polynômes} $P(x)$ à valeurs dans un $\unsure{Gl|GL}(k^n)$, telles que $P(x+y)=P(x)P(y)$.

Les représentations \uline{irréductibles}. Alors le commutant \deleted{\unsure{des}} de l'ens. des $P(x)$ est un corps, \deleted{\unsure{soit donc|comme donc}} soit $D$, et les $P(x)$ prennent leurs valeurs dans le centre $\unsure{Z|\mathfrak z}$ de $D$. \inserted{(Comme $P$ est irréduc., on sait aussitôt que $D=Z$)}

$Z$ est un corps \inserted{\uline{commutatif}} fini sur $k$, donc on a une application $k$-polynômiale de $k$ dans une \deleted{\unsure{corps $Z$}}\inserted{\unsure{extension}} commutative $Z$ de $k$, telle que $P(x+y)=P(x)P(y)$. \deleted{\unsure{Si donc un|Soit donc un}} Je dis que $P(x)=1$ pour tt $x$ ;

\deleted{\unsure{étendant} les corps des scalaires à la clôture algébrique $\hat{k}$ de $k$, on obtient un $\hat{P}(x)$ hom. \inserted{\unsure{polynômial|polynôme}} de $\hat{k}$ dans le groupe multiplicatif des \unsure{élémts} inversibles de la $\hat{k}$-algèbre de dim. finie $A=\hat{k}\otimes_k Z$ \unsure{qui est comm. unitaire|On est ramené aussitôt} \unsure{finie}.}

\deleted{\unsure{On a}} \unsure{on met}

\[\begin{aligned} &\deleted{P(x)=a_n x^n+a_{n-1}x^{n-1}\unsure{\ldots+|+|+\ldots+}a_{\unsure{1|0}}x+1,\quad(\unsure{a_i|a}\in Z,\ x\in k)}\\ &\deleted{\text{avec }a_n\ne 0} \end{aligned}\]

\[\begin{aligned}&\deleted{P(x+y)=a_n x^n+a_n y^n\unsure{+\ldots}}\\ &\deleted{P(x)P(y)=a_n^2 x^n y^n+\unsure{P(x)|\ldots}}\end{aligned}\]

$P$ \unsure{serait} un polynôme de degré $n>0$, $P(x+y)$ \unsure{serait} de degré $n$ et $P(x)P(y)$ de degré $2n$, absurde.

% archive sheet number: 184

\scanpage{186}

Ainsi $P(x)=1$, ce qui implique que $\unsure{Z|\mathfrak z}=k$, donc

\uline{Prop} \deleted{\unsure{toute|Toute}}\inserted{Toute} représentation \inserted{lin.} rationnelle irréductible de $k$ est de dim 1.

\uline{Corollaire} \unsure{Une|Toute} représentation lin. \inserted{rationnelle} de dim finie de $k$ peut se mettre sous forme \deleted{diagonale}\inserted{triangulaire}, avec des $1$ sur la diagonale.

Si on est en caractéristique $p$\inserted{$\ne 0$}, on a \deleted{\unsure{une}} \deleted{\unsure{automorphisme}}\inserted{endomorphisme} \unsure{rationnel|rationnelle|rationnelles} de $k$ \deleted{du type} $F$, $x\unsure{\to|-|\mapsto}x^p$, \deleted{\unsure{qui}} d'où une opération naturelle sur les classes de représentations linéaires, en composant avec $F$ et ses itérés ; ce phénomène est tout à fait distinct de ce qui se passe en car. 0 ; notons que parmi les repr. de dim 2 en car. $p$, on a \unsure{ceux} \deleted{\unsure{qui}} par \deleted{\unsure{rotations type}} translation dans l'espace des polynômes $f(\lambda,\mu)\unsure{:|=|(x)=}\lambda X^{p^h}+\mu$ ($h$ fixé, $X$ l'indéterminée) :

\[\begin{aligned}\tau_y f(\lambda,\mu)\unsure{(X)|\in X)|(x)} &= \lambda(X+y)^{p^h}+\mu\\ &=\lambda X^{p^h}+\mu+\lambda y^{p^h}\unsure{,=|=}f(\lambda,\mu+\lambda y^{p^h}),\quad \text{i.e.}\end{aligned}\]

la représentation $E_h(y)=\begin{pmatrix}1&y^{p^h}\\0&1\end{pmatrix}$ ;

les représentations, pour \deleted{\unsure{des}} \deleted{\unsure{variables}} $h$ distincts, sont inéquivalentes deux à deux, d'où une infinité de classes de repr. lin. de dim 2. En car. 0, il n'y en a que deux, au contraire.

% archive sheet number: 185

\scanpage{187}

Supposons maintenant la caractéristique nulle. Alors les représentations \inserted{lin.} rationnelles du \uline{groupe} \uline{abélien} $k$ correspondent à \uline{certaines} représentations lin. de l'\uline{algèbre de Lie abélienne} $k$ ; il résulte de la \deleted{\unsure{pr.}} prop. 1 qu'il s'agit exactement des représentations \uline{nilpotentes} de l'algèbre $k$, dont la classification est celle des $k[X]$-modules de type fini $X$-primaires. Ainsi :

\noindent\begin{tabular}{@{}|p{.92\linewidth}@{}}\uline{Prop. 2} La classification des repr. linéaires du groupe alg. abélien $k$ est, via les représentations infinitésimales associées, équivalente à la classification des $k[X]$-modules de type fini $(X)$-primaires.\end{tabular}

% slanted
\margin{$k$ de car $0$}

Les $k[X]$-modules $(X)$-primaires \uline{indécomposables} sont ceux qui sont du type $k[X]/X^{n}$ $(1\leq n<+\infty)$. On \deleted{\unsure{du}} voit aussitôt que le dual \inserted{sur $k$} du $k[X]$ module $E_n$ est isomorphe à $E_n$.

% slanted
\margin{La classification qu'on a donné des \uline{classes d'\unsure{opérateurs|invariants}} \uline{nilpotents} \deleted{\illegible{1}} dépend \uline{aussi} de la caractéristique.}

Ainsi

\uline{Corollaire} L'anneau \inserted{$\unsure{A|\mathcal A}(k)$} des classes de représentations du groupe alg. $k$ a pour base canonique

% archive sheet number: 186

\scanpage{188}

\noindent\begin{tabular}{@{}|p{.92\linewidth}@{}}les classes $E_n$ des modules $k[\unsure{x|X}]/x^n$, $n\unsure{:|=}1,2,\ldots$. L'application $\unsure{v|\vee}$ dans $A$ est réduite à l'identité. On a $d(E_n)=n$, $E_1$ est l'unité.\end{tabular}

Structure multiplicative. Posons $\xi=E_2$, classe de $k[x]/(X^{\unsure{2}})$. Alors $\xi^n$ est la classe de $F_n=k[X_1,\ldots,X_n]/(X_1^2,\ldots X_n^2)$, muni de l'opérateur nilpotent $T$ de multiplication par la classe $\dot S$ de $S=X_1+\ldots+X_n$. \deleted{Les} Le \unsure{$n$-ième} itéré \inserted{$T^h$} de $T$ est la multiplication par

\[ \deleted{\dot S^h=h!\sum_{\substack{\unsure{H\subset[1,n]}\\\unsure{\lvert H\rvert=h|\#H}}}}\qquad \dot S^h=k!\sum_{1\leq i_1<\cdots<i_{\unsure{k|h}}\leq n}\dot X_{i_1}\ldots\dot X_{i_{\unsure{h|k}}}. \]

En particulier $T^n$ est la multiplication par $\dot S^n=n!\dot X_1\ldots\dot X_n$, et $T^{n+1}$ est nul.

On en conclut aussitôt que $T^n F_n$ est de dim 1 (base formée de $\dot X_1\dot X_2\ldots\dot X_n$) et $T^{n+1}(F_n)=0$ ; d'où résulte facilement que \deleted{\unsure{la multiplicité}} $\xi^n$ est combinaison linéaire des $E_1,\ldots E_n$ à coefficients $\in Z^+$, celui de $E_n$ étant 1 :

\margin{\unsure{3}}

% archive sheet number: 187

\scanpage{189}

\[\xi^n=L_n(E_1,\ldots,E_n)=\sum_{i=1}^{n}c_{ni}E_i\qquad\text{\uuline{$c_{nn}=1$}}\]

On en conclut la

\noindent\begin{tabular}{@{}|p{.92\linewidth}@{}}\uline{Proposition 3} $\unsure{A|\mathcal A}(k)$ est l'anneau de polynômes $Z[\xi]$ sur $Z$ engendré par $\xi=E_2$ $(d(\xi)=2)$\end{tabular}

N. B. Le calcul explicite donne

\diagram{\(\begin{array}{l@{\qquad\qquad}l}\xi^0=E_1&E_1=1\\\xi^1=E_2&E_2=\xi\\\xi^2=E_1+E_3&E_3=\xi^2-1\\\xi^3=2E_2+E_4&E_4=\xi^3-2\xi\\\xi^4=2E_1+3E_3+E_5&E_5=\xi^4-3\xi^2+5\\\ldots&\ldots\end{array}\)}

(Toutes les structures de $A(k)$ sont bien précisées dans \unsure{ce|un} th. de structure, \uline{sauf} la structure d'ordre, qui est un peu cachée !)

% archive sheet number: 188

\scanpage{190}

$\underline{\overline{\mathrm{II}}}$\quad $\unsure{A|\mathcal A}(k^*)$. \deleted{\unsure{Le groupe}} \unsure{a)|1)|\textsuperscript{$\circ$})} Représentations irréductibles.

Se \unsure{fait|tout|font} via une représentation \unsure{rationelle|rationnelle} de $k^*$ dans $L^*$, où $L$ est un corps extension finie de $k$. $P(\lambda)=\sum_{-n}^{n}\lambda^i c_i$ $(c_i\in L)$

on a $P(\lambda\lambda')=P(\lambda)P(\lambda')$, i.e.

\[\sum_{\deleted{\unsure{i=-\infty|i}}}^{\unsure{\infty|n}}\lambda^i{\lambda'}^i c_i=\sum_{ij\deleted{\unsure{-\infty}}}^{\unsure{\infty|n}}\lambda^i{\lambda'}^j c_i c_j\unsure{=|,}\ ,\quad \text{et comme}\]

c'est une identité en $\lambda,\lambda'$ on doit avoir $c_i=c_i^2$, $c_i c_j=0$ si $i\ne j$, donc il ne peut y avoir qu'\deleted{\unsure{un des}} \uline{un} des $c_i$ qui soit $\ne 0$, soit $c_n$, et $\unsure{c=c^2|c_n=c_n^2}$ exige $c_n=1$. Donc $P(\lambda)=\lambda^n$ ; l'irréductibilité exige $L\unsure{:|=}k$. Ainsi

\margin{$k$ \uline{infini} (?)}

\noindent\begin{tabular}{@{}|p{.92\linewidth}@{}}\uline{Proposition} Les seules repr. \deleted{\unsure{idél.|irréd.}} irréductibles de $k^*$ ($k$ corps \uline{infini}) sont les repr. de dim 1, équivalentes aux repr. $\deleted{\unsure{\lambda\to\lambda^n}}\ \sigma_n(\lambda)t=\lambda^nt$.\end{tabular}

\uline{Corollaire} $\sigma_n\sigma_m=\sigma_{n+m}$, $\sigma_n^{\unsure{\prime|\vee}}=\sigma_{-n}$

Plus généralement, considérons une repr. lin. \inserted{\unsure{rationelle|rationnelle}} \uline{quelconque} de ${\unsure{\lambda|k}}^*$ dans un vectoriel $V$. On peut la mettre sous la forme

% archive sheet number: 189

\scanpage{191}

\[P(\lambda)=\sum_0^n c_i\lambda^i\qquad(c_i\in L(\unsure{V|W}))\]

On voit comme ci-dessus que

\[c_i c_j=0\quad\text{si}\quad i\ne j\qquad c_i^2=c_i\]

De plus $\sum c_i=P(1)=1$, donc les $c_i$ sont des projecteurs deux à deux orthogonaux, ainsi

\uline{Proposition} Toute représentation \inserted{lin. rat.} de $C^*$ est complètement réductible ($k$ corps infini quelconque).

\uline{Corollaire} $\unsure{A|\mathcal A}(k^*)\simeq Z(\deleted{\unsure{F}}\unsure{\Gamma|F|Z})$, $\deleted{\unsure{G|F}}\Gamma$ groupe des \unsure{entiers|caractères} rationnels ; dans cet isomorphisme, $\sigma_n\longleftrightarrow$ élément de $Z(\Gamma)$ corresp. \deleted{\unsure{en}} à $n\in\Gamma$.

L'isomorphisme est compatible avec l'augmentation, \unsure{relation|notion} d'ordre, \inserted{$\unsure{v|\vee}$.}

Différentielle de \inserted{la repr.} de dim \unsure{1|$\lambda^n$} : $t\unsure{\to|-}nt$.

\uline{Généralisation aux ``tores'' $G$ isomorphes aux $\unsure{C^{*n}|(k^*)^n}$.}

Si $\Gamma$ est le groupe des repr. \unsure{rati.|rat.|r.t.} de $G$ dans $k^*$, $\Gamma$ est $\simeq Z^n$, et $A(G)\simeq Z(\Gamma)$.

% archive sheet number: 190

\scanpage{192}

$\underline{\overline{\mathrm{III}}}$\quad $G=\underset{\substack{\parallel\\F}}{k}\mathbin{.}\underset{\substack{\parallel\\H}}{k^*}$ produit semi-direct, $k^*$ opérant dans $k$ par $\sigma_p(\lambda)t=\lambda^pt$, ($n$ entier rationnel \uline{non nul}).

On suppose $k$ de car. 0. L'algèbre de Lie \deleted{\unsure{de}} $\unsure{\mathfrak g|\mathscr G}$ de $G$ est extension de $\unsure{\mathscr F|\mathfrak F|F}\simeq k$ par $\unsure{\mathscr H|\mathfrak H|H}\simeq k$, $\mathscr H$ opérant sur $\mathscr F$ \deleted{par} ainsi ($X$ et $Y$ sont les \unsure{éléments-base|éléments-bases} canoniques de $\mathscr F$ et $\mathscr H$) :

\[[Y,X]=pX\tag*{(1)}\]

Considérons une représentation linéaire indécomposable \unsure{rationnelle|rationelle} $\varphi$ de $\deleted{\unsure{G_2(\mathscr F)}}G$, soit $\varphi'$ sa différentielle. On a (décomposition en sous-espaces propres \unsure{relt|rel\textsuperscript{ts}} à $H\subset G$)

\[V=\sum V_n\]

($Y.a=\deleted{p}na$ si $a\in V_n$). La condition (1) montre que les composantes $X_{\deleted{\unsure{ij|i}}ji}\in L(V_i,V_j)$ de $X$ satisfont :

\[(j-i)X_{ji}=pX_{ji}\]

d'où $X_{ji}=0$ sauf si $j-i=p$.

% archive sheet number: 191

\scanpage{193}

Si tous les $X_{ji}$ sont nuls, alors ($V$ étant indécomposable) il n'y a qu'un seul $V_n$ non nul, qui est de dimension un. On désignera la repr. ainsi obtenue (définie par la repr. $\sigma_n$ de $H=C^*$ en composant avec $G\longrightarrow H$) par \uline{$\sigma_n$}. S'il existe un $X_{\beta\alpha}$ non nul, alors \inserted{$\beta-\alpha=p$, donc} on a $V=V'+V^{\unsure{\prime\prime\prime|\prime\prime}}$,

où $V'=\sum_{\substack{n\equiv\alpha\\(p)}}V_n$, $V''=\sum_{\substack{n\not\equiv\alpha\\(p)}}V_n$

sont stables sous $\unsure{\mathfrak g|\mathscr G}$, et $V'\ne0$, donc $V'=V$. \deleted{Soit} On a donc, si \deleted{\unsure{$p$ est le plus}} \inserted{$V_n,V_{n'}\ne0$ : $n'-n\equiv0\ (p)$.}

\inserted{Soit $n$ le plus} petit entier tel que $V_{\unsure{n'|n}}\ne0$, alors les $V_m$ éventuellement non nuls \deleted{\unsure{du}} sont de la forme $V_{n+pk}$ $k=0,1,\ldots$. Soit $k$ le plus grand entier tel que $V_{n+pk}\ne0$, \deleted{\unsure{alors}} on a $k\geq1$. Soit $u_i=X_{n+p(i+1),n+pi}$,

\deleted{\unsure{on a} \unsure{$a=0$} $V$} $E_i=V_{n+pi}$, on a donc une suite d'applications

\diagram{\begin{tikzcd}[ampersand replacement=\&,column sep=small] E_0\arrow[r,"u_0"] \& E_1\arrow[r,"u_1"] \& E_2 \& \ldots\arrow[r,"u_{k-1}"] \& E_k. \end{tikzcd}}

\uline{Lemme 1}. Les $u_i$ sont \uline{injectifs}. \inserted{$(0\leq i\leq k-1)$} On va le prouver par récurrence sur $i$, c'est trivial pour $i=-1$ : $0=E_{-1}\longrightarrow E_0$. Supposons

% archive sheet number: 192

\scanpage{194}

le fait jusqu'à $u_{\ell-1}$ \deleted{\unsure{$(\ell\leq k)$}}\inserted{$(0\leq\ell\leq k-1)$} prouvons le pour $u_\ell$.

\deleted{Soit} Soit $x$ un élément de $\operatorname{Ker}u_\ell$, \deleted{\unsure{on va montrer}, \unsure{prouvons} qu'il appartient à tout $E_{\ell'}$ $(\ell'\leq\ell)$ (identifié à un sous-espace de $E_\ell$), \unsure{par|$+0$} \unsure{récurrence} et on aura fini. \unsure{On prend}} \inserted{s'il n'était} pas nul il existerait un \deleted{premier} \deleted{\unsure{dernier}} \inserted{\unsure{plus haut|plus grand}} $\ell'$ \deleted{\unsure{$(0\leq\ell'<\ell)$|$(0\leq\ell'\leq\ell)$}}\inserted{$(-1\leq\ell'<\ell)$}

tel que $x\notin E_{\ell'}$ -- $E_{\ell'}$ identifié à un sous-espace de $E_\ell$. Soit $D$ la droite engendrée par $x$, \deleted{soit} $(D\subset E_{\ell'+1})$ et soit $\deleted{\unsure{E|E_j}}F$ un hyperplan de $E_\ell$ \deleted{\unsure{de}} contenant $E_{\ell'}$ et suppl. de $D$ ; soit pour $0\leq j\leq k$,

\margin{$E_{\ell'}\quad E_{\ell'+1}\ldots E_\ell$}

\[\deleted{E'_j\unsure{=}\left\{\begin{array}{ll}E_j&\text{si }j\leq\ell'\\D&\text{si}\end{array}\right.}\]

\[E'_j=E_j\deleted{\unsure{\cap F}}\quad\text{si}\quad j\leq\ell'\quad\text{ou}\quad\deleted{\unsure{E_j=E_j|E'_j=E_j}}\inserted{j>\ell},\quad\deleted{\unsure{j>\ell}}\inserted{E'_j=E_j\cap F\quad\text{si }\ell'<j\leq\ell}\]

\[ E''_j=0\quad\text{si}\quad j\leq\ell'\ \text{ou}\ j>\ell,\qquad E''_j=D\quad\text{si }\ell'<j\leq\ell \]

Alors $E_j=E'_j\oplus E''_j$ pour tout $j$, \deleted{\unsure{les $E'_j$ et}} \inserted{et $\sum E'_j$ et} $\sum E''_j$ sont stables sous $\deleted{\unsure{D}}\unsure{\mathfrak g|\mathscr G}$ et $\ne0$, donc $\sum E_i$ ne serait pas indécomposable.

\uline{Lemme 2} Les $u_i$ $(0\leq i\leq k-1)$ sont \uline{surjectifs}. Se voit soit directement de façon analogue, soit en \deleted{\unsure{tenant}} appliquant le lemme 1 à la repr. contragrédiente.

% archive sheet number: 193

\scanpage{195}

Conclusion : les $u_{\deleted{\unsure{j}}i}$ \inserted{$(0\leq i\leq k-1)$} sont des bijections, les $E_i$ sont de dim. 1.

% slanted
La représentation de $\unsure{\mathfrak g|\mathscr G}$ ou de $G$ qu'on vient de construire est notée $e_{n,k+1}$ $(k+1\geq2)$ \inserted{\unsure{Or|Car} la repr. ainsi construite est indécomposable, car elle induit déjà sur $F$ la repr. indécomposable notée $e_{k+1}$ plus haut} Si $k=0$, il est naturel de poser $e_{n,1}=\sigma_n$.

Ainsi

\noindent\begin{tabular}{@{}|p{.92\linewidth}@{}}\uline{Proposition} Les représentations indécomposables de $G$ sont \deleted{les}\inserted{équivalentes aux} représentations $e_{n,k}$ \deleted{\unsure{$(k\geq1)$}} $(n\in Z,\ k\in Z,\ k\geq1)$ qu'on vient de construire : \deleted{\unsure{$\sigma$}} $e_{n,1}$ est la représentation $\sigma_n$ de dim 1 induite par la représentation de $=$ nom de $H=C^*$, $e_{n,k}=\sigma_n\otimes e_{0,k}$, où $e_{0\,k}$ est la représentation \deleted{\unsure{à val}} dans l'espace $k[Z]/Z^k$ définie ainsi : $X$ \unsure{définit|désigne|étant} l'opérateur de multiplication par $Z$ $\bmod (Z^k)$, et

\[Y.Z^\alpha=\alpha p Z^\alpha\quad(Z^k).\]

\inserted{les plus bas et plus hauts poids de $e_{n,k}$ sont $n$ et $n+p(k-1)$, donc} les représentations $e_{n,k}=\sigma_n\otimes e_{0\,k}$ sont deux à deux inéquivalentes. $e_{n,k}$ est de degré $k$ ; $\overset{\vee}{e}_{n,k}=e_{-n-p(k-1),k}$.\end{tabular}

% archive sheet number: 194

\scanpage{196}

Nous allons déterminer la structure multiplicative de $A(G)$. Les \unsure{hom.} \diagram{\begin{tikzcd}[ampersand replacement=\&] F\arrow[r,dash,"i"] \& G\arrow[r,dash,"j"] \& H \end{tikzcd}} définissent des homomorphismes

\diagram{\begin{tikzcd}[ampersand replacement=\&] A(H)\arrow[r,"j^*"] \& A(G)\arrow[r,"i^*"] \& A(F) \end{tikzcd}}

$A(H)$ est l'algèbre \inserted{sur $Z$} engendrée par les $\sigma_n$ $(n\in Z)$ soumis aux relations $\sigma_n\sigma_{n'}=\sigma_{n\unsure{-|+}n'},\ \sigma_0=1$. Il est plongé dans $A(G)$ par $j^*$, par

\[ j^*(\sigma_n)=e_{n\,1}. \]

$A(F)$ est \deleted{\unsure{l'algèbre en}} une algèbre ayant une base formée d'éléments $e_n$ $(n\geq1)$ \deleted{\unsure{et}} $e_1=1$. D'ailleurs $A(F)=Z[e_2]$. On a

\[ i^*(e_{nk})=e_k \]

Identifions $\Lambda=A(H)$ à une sous-algèbre de $A(G)$. Posons $\xi=e_{0\,2}$. On se propose de prouver que

\[ A(G)=\Lambda[\xi]\qquad\text{(alg. des polynômes \unsure{en $\xi$,|à une $\xi$,} \unsure{cff.|coeff.} \unsure{dans $\Lambda$)|de $\Lambda$)}} \]

i.e. que les $\xi^q$ $(q=0,1,\ldots)$ forment une \uline{base} de $A(G)$ sur $\Lambda$. \deleted{\unsure{Pour}}\inserted{Comme} les $e_{0\,1},e_{0\,2},e_{0\deleted{\unsure{j}}3}\ldots$ forment une telle base, il \deleted{\unsure{on}}, il suffit d'établir qu'on a des formules

\[ \xi^q=\sum_{i=1}^{q+1}c_{qi}e_{0\,i}\qquad(c_{qi}\in\Lambda,\ c_{q,q+1}=1)\tag*{(1)} \]

% archive sheet number: 195

\scanpage{197}

Or à priori on a

\[ \xi^{n}=\sum_i c_{\unsure{q|p|n}i}\,e_{0\,i} \]

les $c_{qi}\in\Lambda^+$, d'où en \unsure{transformant} par $i^*$ :

\[ e_2^n=\sum_i i^*(c_{qi})\,e_i \]

\deleted{Or \unsure{on a} $i^*(c_{p}$}

Or on a

\[ i^*(\lambda\deleted{\unsure{p}})=d(\lambda)\qquad\lambda\in\Lambda \]

$d$ étant l'augmentation de $\Lambda$, d'où

\[ e_2^n=\sum_i d(c_{qi})\,e_{\unsure{i}} \]

Or on sait déjà qu'on a $e_2^q=\sum_1^{q+1}a_{qi}\,e_i$, les $a_{qi}$ étant des entiers $\geq0$ et \deleted{\unsure{tels}} $a_{q,q+1}=1$. Donc on a

\[ \varepsilon(c_{\deleted{\unsure{n|p}}qi})=0\quad\text{si}\quad i>q+1,\qquad \varepsilon(c_{q,q+1})=1 \]

d'où, en raison des $c_{qi}\in\Lambda^+$,

\[ c_{qi}=0\quad\text{si}\quad i>q+1,\qquad c_{q,q+1}\ \text{de la forme}\ \sigma_\alpha \]

($\alpha$ dépendant de $n$ \deleted{\unsure{$q$}}). D'où les formules (1), mais on doit encore déterminer \deleted{\unsure{$=$}} \deleted{\unsure{$c_{p,p+1}=\sigma_\alpha$}.} \deleted{\illegible{3}} $c_{p,p+1}=\sigma_\alpha$

Or le plus haut poids dans la repr. $\sigma_\alpha e_{0\,p+1}=e_{\alpha,p+1}$ est $\alpha+pn$, \deleted{et} le plus

% archive sheet number: 196

\scanpage{198}

haut poids pouvant figurer dans $\xi^n$ est $n$ fois le plus haut poids figurant dans $\xi$, soit $pn$. On a donc $pn+\alpha\leq pn$, d'où $\alpha\leq0$.

D'autre part, les poids intervenant dans $\xi$ sont $0$ et $p$, donc positifs, donc ceux dans $\xi^n$ sont positifs ; il faut donc qu'il \deleted{\unsure{$\alpha$}} en \deleted{soit} de $=$ des poids intervenant dans $\sigma_\alpha e_{0,p+1}$, d'où $\alpha\geq0$. D'où $\alpha=0$, i.e. $\sigma_\alpha=1$, c.q.f.d.

Ainsi :

% vertical rule at the left of the proposition
\noindent\begin{tabular}{@{}|p{.92\linewidth}@{}}\uline{Proposition} $A(G)$ est l'algèbre de polynômes sur $\Lambda=Z(\Gamma)$ $(\Gamma=Z)$ engendrée par $\xi=e_{0\,2}$ ; l'augmentation de $\Lambda$ est son augmentation usuelle, et $d(\xi)=\deleted{p}2$. Sur $\Lambda$, l'involution $\vee$ est l'involution usuelle, et on a
\[ \check\xi=e_{-p,2}=\sigma_{-p}.\xi. \]
\end{tabular}

% archive sheet number: 197
